%!TeX program = lualatex
\renewcommand{\thelektion}{L10 \textperiodcentered{} Aufbau eines Computers I}

\begin{frame}\lessontitle{Lektion 10}{Aufbau eines Computers I}{Von-Neumann-Architektur \& die CPU im Detail \textperiodcentered{} Skript, Kapitel 5}\end{frame}

\begin{frame}\FDUtitle{Vom Gatter zum Computer}
\begin{center}
Wir kennen jetzt Gatter und den Halbaddierer.
\vspace{4mm}

Wie setzt man daraus einen \textbf{vollständigen Computer} zusammen?
\end{center}
\begin{center}
\begin{tikzpicture}[node distance=1.5cm]
    \node[draw, fill=FDUblue!15, minimum width=2.5cm, minimum height=1cm, rounded corners] (ein) {Eingabe};
    \node[draw, fill=FDUgreen!35, minimum width=2.5cm, minimum height=1cm, rounded corners, right=of ein] (ver) {Verarbeitung};
    \node[draw, fill=FDUteal!25, minimum width=2.5cm, minimum height=1cm, rounded corners, right=of ver] (aus) {Ausgabe};
    \draw[->, thick] (ein) -- (ver);
    \draw[->, thick] (ver) -- (aus);
\end{tikzpicture}
\end{center}
\end{frame}

{\setbeamertemplate{footline}{\darkfootline}
\begin{frame}\sectionframe[FDUblue]{Die Von-Neumann-Architektur}{1945, John von Neumann}
\end{frame}}

\begin{frame}\FDUtitle{Das entscheidende Merkmal}
\begin{center}
\textbf{Programm und Daten} liegen im \textbf{gleichen} Speicher.
\vspace{4mm}

Das Programm ist selbst nur ein Stück Daten -- also eine Bitfolge, die jederzeit verändert oder ausgetauscht werden kann.
\end{center}
\end{frame}

\begin{frame}\FDUtitle{Die Hauptkomponenten}
\begin{center}
\begin{tikzpicture}[
    block/.style={draw, minimum width=2.4cm, minimum height=1.0cm, text centered, rounded corners=3pt, font=\small},
    every path/.style={thick}
]
    % CPU box
    \node[draw, dashed, minimum width=8.8cm, minimum height=2.4cm, rounded corners=4pt, label={north:{\textbf{CPU (Prozessor)}}}] (cpu) at (0,0) {};

    \node[block, fill=FDUgreen!25] (alu) at (-2.8, 0) {ALU};
    \node[block, fill=Gray!20, minimum width=2.0cm] (reg) at (0, 0) {Register};
    \node[block, fill=FDUblue!20] (cw) at ( 2.8, 0) {Steuerwerk};

    \draw[<->] (alu) -- (reg);
    \draw[<->] (reg) -- (cw);

    % RAM
    \node[block, fill=FDUteal!25, minimum width=4.0cm] (ram) at (0, -2.5) {Arbeitsspeicher (RAM)};

    % Systembus
    \draw[<->, line width=1.5pt] (cpu.south) -- (ram.north) node[midway, right=3pt, font=\small] {Systembus};

    % Massenspeicher & IO
    \node[block, fill=Orchid!25, minimum width=2.8cm] (stor) at (-3.7, -2.5) {Massenspeicher};
    \node[block, fill=Salmon!30, minimum width=2.8cm] (io) at ( 3.7, -2.5) {Ein-/Ausgabe};

    \draw[<->] (stor) -- (ram);
    \draw[<->] (ram) -- (io);
\end{tikzpicture}

\end{center}
\end{frame}

\begin{frame}\FDUtitle{Warum ist das riskant?}
\begin{center}
Was könnte passieren, wenn ein fehlerhaftes Programm versehentlich seinen eigenen Programmcode überschreibt?
\end{center}
\only<2->{
\begin{center}\color{FDUteal}\bfseries Unvorhersehbares Verhalten oder Absturz -- deshalb übernimmt später das Betriebssystem den \enquote{Speicherschutz} (Lektion 12).\end{center}
}
\end{frame}

{\setbeamertemplate{footline}{\darkfootline}
\begin{frame}\sectionframe[FDUgreen]{Die CPU im Detail}{ALU, Steuerwerk, Register}
\end{frame}}

\begin{frame}\FDUtitle{ALU -- Arithmetisch-logische Einheit}
\begin{itemize}\setlength\itemsep{7pt}
  \item Das \enquote{Rechenzentrum} der CPU
  \item Arithmetische Operationen (Addition, Subtraktion, \ldots)
  \item Logische Operationen (AND, OR, Vergleiche, \ldots)
  \item Intern: eine grosse Anzahl kombinierter Logikgatter (z.\,B.\ Volladdierer)
\end{itemize}
\end{frame}

\begin{frame}\FDUtitle{Steuerwerk \& Register}
\begin{columns}[T,onlytextwidth]
\begin{column}{.46\textwidth}
\large\bfseries\color{FDUteal}Steuerwerk
\vspace{2mm}
\normalsize
\begin{itemize}\setlength\itemsep{5pt}
  \item Steuert den Ablauf aller Operationen
  \item Liest Befehle, interpretiert sie
  \item Gibt Steuersignale an ALU \& Co.
\end{itemize}
\end{column}
\begin{column}{.46\textwidth}
\large\bfseries\color{FDUblue}Register
\vspace{2mm}
\normalsize
\begin{itemize}\setlength\itemsep{5pt}
  \item Sehr schnelle, winzige Zwischenspeicher
  \item Halten die aktuellen Rechenwerte
  \item Typische Grösse: 64 Bit (= 8 Byte)
\end{itemize}
\end{column}
\end{columns}
\end{frame}

\begin{frame}{Auftrag}{Kapitel 5: Von-Neumann-Architektur \& CPU}
\aufgabebox{\textbf{Aufgabe 5.1}: Überlegen Sie, warum es problematisch ist, wenn Programm und Daten im selben Speicher liegen.}
\end{frame}
